
\noindent \uline{When the two sides of} (7.3.1) \uline{are regarded as modules over} @@M0@@ \uline{by means of the ring homomorphisms} @@M1@@ \uline{and} @@M2@@ (IV 2.7 b) \uline{and the module structures on} @@M3@@ \uline{and} @@M4@@ \uline{over} @@M5@@ \uline{and} @@M6@@ \uline{respectively} (IV 2.10), \uline{the homomorphism} @@M7@@ \uline{is a homomorphism of} @@M8@@-\uline{modules}.

The last assertion follows from the fact that @@M9@@ is a homomorphism of @@M10@@-modules, as is the epimorphism @@M11@@, by (IV 2.12.2).

We shall review other properties of the specialization morphism.

\sgaStableTarget{sga6.e10.sec.7.4}\subsection*{7.4. Compatibility with filtrations.}
Endow the @@M12@@ with the filtration @@M13@@ by the dimension of supports (1.1.1). Then, when @@M14@@ is of finite type, more generally when @@M15@@ is catenary, \emph{the specialization homomorphism is compatible with these filtrations}, and consequently induces a homomorphism (also called the specialization homomorphism):
\sgaStableTarget{sga6.e10.eq.7-4-1}@@M16@@

To see this, by 1.1.4 one is reduced to proving that, when @@M17@@ is an integral subscheme of dimension @@M18@@ of @@M19@@, @@M20@@ has filtration @@M21@@. Now let @@M22@@ be the schematic closure of @@M23@@ in @@M24@@, so that @@M25@@ is flat over @@M26@@ (EGA IV 2.8.1). It follows readily that @@M27@@, where the two @@M28@@ are taken in @@M29@@ and @@M30@@ respectively, whence @@M31@@. One checks readily (using the fact that @@M32@@ is catenary, hence @@M33@@ is catenary) that @@M34@@; when @@M35@@ is of finite type, this follows at once, by reduction to the case of integral @@M36@@, from EGA IV 7.1.13. This gives the asserted conclusion.
